% Chapter 2, Section 3 _Linear Algebra_ Jim Hefferon
%  http://joshua.smcvt.edu/linearalgebra
%  2001-Jun-10
\section{基与维数}
\index{basis|(}
前一节的结尾指出：张成集在它线性无关时是最小的，
而线性无关集在它张成整个空间时是最大的。
所以“最小张成集”与“最大无关集”这两个概念恰好重合。
本节将给这个想法命名，并研究它的性质。





\subsection{基}
\begin{definition} \label{def:basis}
%<*df:basis>
向量空间的\definend{基}\index{basis!definition}\index{vector space!basis}
是由线性无关且张成该空间的向量组成的序列。
%</df:basis>
\end{definition}

%<*BasisNotation>
由于基是一个序列，
也就是说，
两组基即使所含元素相同而次序不同，
也算不同的基，
% the 
% order of the elements is significant,
所以我们用尖括号来记它：
\( \sequence{\vec{\beta}_1,\vec{\beta}_2,\ldots} \).\appendrefs{sequences}\@ %\spacefactor=1000{}
%</BasisNotation>
（若由序列的元素构成的重集是无关的，则称该序列是线性无关的。
类似地，
若序列的元素构成的集合张成该空间，
则称该序列张成该空间。）

\begin{example}  \label{ex:FstBasisRTwo}
这是 \( \Re^2 \) 的一个基。
\begin{equation*}
  \sequence{ \colvec{2 \\ 4},\colvec{1 \\ 1} }
\end{equation*}
它是线性无关的
\begin{equation*}
  c_1\colvec{2 \\ 4}+c_2\colvec{1 \\ 1}=\colvec{0 \\ 0}
  \quad\implies\quad
  \begin{linsys}{2}
     2c_1  &+  &1c_2  &=  &0  \\
     4c_1  &+  &1c_2  &=  &0  
   \end{linsys}
  \quad\implies\quad
  c_1=c_2=0
\end{equation*}
并且它张成 \( \Re^2 \)。
\begin{equation*}
  \begin{linsys}{2}
    2c_1  &+  &1c_2  &=  &x  \\
    4c_1  &+  &1c_2  &=  &y  
  \end{linsys}
  \quad\implies\quad
  c_2=2x-y\text{\ 而\ } c_1=(y-x)/2
\end{equation*}
\end{example}

\begin{example}
\( \Re^2 \) 的这个基与前面那个不同
\begin{equation*}
  \sequence{\colvec[r]{1 \\ 1},\colvec[r]{2 \\ 4}}
\end{equation*}
这是因为其中元素的次序不同。
验证它是一个基，做法与前面的例题完全相同。
\end{example}

\begin{example}
空间 \( \Re^2 \) 有许多基。
另一个是这样的。
\begin{equation*}
  \sequence{ \colvec[r]{1 \\ 0},\colvec[r]{0 \\ 1} }
\end{equation*}
验证很容易。
\end{example}

\begin{definition} \label{df:StandardBasis}
%<*df:StandardBasis>
对任意 \( \Re^n \)
\begin{equation*}
   \stdbasis_n=\sequence{
     \colvec[r]{1 \\ 0 \\ \vdotswithin{0} \\ 0},
     \colvec[r]{0 \\ 1 \\ \vdotswithin{0} \\ 0},
     \dots,\,
     \colvec[r]{0 \\ 0 \\ \vdotswithin{0} \\ 1}}
\end{equation*}
是\definend{标准}\index{standard basis}%
\index{basis!standard}基（也称为\definend{自然}基）。
我们把这些向量记为 \( \vec{e}_1,\dots,\vec{e}_n \)。
%</df:StandardBasis>
\end{definition}

\noindent 
微积分教材把 $\Re^2$ 的标准基向量记作
\( \vec{\imath} \) 和 \( \vec{\jmath} \)，而不是 $\vec{e}_1$
和 $\vec{e}_2$；并把 \( \Re^3 \) 的标准基向量记作
\( \vec{\imath} \)、\( \vec{\jmath} \) 和 \( \vec{k} \)，
而不是 $\vec{e}_1$、$\vec{e}_2$ 和 $\vec{e}_3$。
注意 \( \vec{e}_1 \) 在讨论 \( \Re^3 \) 时的含义
与在讨论 \( \Re^2 \) 时的含义不同。

\begin{example}  \label{ex:BasisForCosPlusSin}
考虑由实变量 $\theta$ 的函数组成的空间
\( \set{a\cdot\cos\theta+b\cdot\sin\theta\suchthat a,b\in\Re} \)。
它的一个自然基是
$ \sequence{\cos\theta, \sin\theta}=\sequence{1\cdot\cos\theta+0\cdot\sin\theta,
             0\cdot\cos\theta+1\cdot\sin\theta}
    $。
这个空间一个更一般的基是
\( \sequence{\cos\theta-\sin\theta,
             2\cos\theta+3\sin\theta} \)。
这两个都是基的验证见
\nearbyexercise{exer:VerifBasesCosPlusSin}。
\end{example}

\begin{example}
三次多项式向量空间 \( \polyspace_3 \) 的一个自然基是
\( \sequence{1,x,x^2,x^3} \)。
这个空间的另外两个基是 \( \sequence{x^3,3x^2,6x,6} \)
和 \( \sequence{1,1+x,1+x+x^2,1+x+x^2+x^3} \)。
验证它们各自都线性无关且张成整个空间，这很容易。
\end{example}

\begin{example}
平凡空间\index{trivial space}\index{vector space!trivial} 
$\set{\zero}$ 只有一个基，即空的基
\( \sequence{} \)。
\end{example}

\begin{example}
由有限次多项式构成的空间有一个含有无穷多个元素的基
\( \sequence{1,x,x^2,\ldots} \)。
\end{example}

\begin{example}
基我们以前就见过。
在第一章中，对于诸如这样的齐次方程组
\begin{equation*}
   \begin{linsys}{4}
      x  &+  &y  &   &   &-   &w   &=  &0  \\
         &   &   &   &z  &+   &w   &=  &0  
   \end{linsys}
\end{equation*}
我们通过参数化来描述其解集。
\begin{equation*}
  \set{\colvec[r]{-1 \\ 1 \\ 0 \\ 0}y
       +\colvec[r]{1 \\ 0 \\ -1 \\ 1}w
       \suchthat y,w\in\Re }
\end{equation*}
于是，解的向量空间就是
一个二元集合的张成。
这个由两个向量组成的集合也是线性无关的，这容易验证。
因此，解集是 \( \Re^4 \) 的一个子空间，以上面这两个向量
组成它的一个基。
\end{example}

\begin{example}
参数化不仅能求出齐次方程组解集的基，
也能求出其他向量空间的基。
为求 $\matspace_{\nbyn{2}}$ 的这个子空间的基
\begin{equation*}
  \set{\begin{mat}
         a  &b  \\
         c  &0
       \end{mat} \suchthat a+b-2c=0}
\end{equation*}
我们把条件改写为 $a=-b+2c$。
\begin{equation*}
  \set{\begin{mat}
         -b+2c  &b  \\
          c     &0
       \end{mat} \suchthat b,c \in \Re}
  =\set{b\begin{mat}[r]
         -1  &1  \\
          0  &0
       \end{mat}+
       c\begin{mat}[r]
          2  &0  \\
          1  &0
       \end{mat} \suchthat b,c \in \Re}
\end{equation*}
于是，下面是一个自然的候选基。
\begin{equation*}
  \sequence{\begin{mat}[r]
         -1  &1  \\
          0  &0
       \end{mat},
       \begin{mat}[r]
          2  &0  \\
          1  &0
       \end{mat} }
\end{equation*}
上面的推导表明它张成该空间。
它的线性无关性也容易证明。
\end{example}

再考虑 \nearbyexample{ex:FstBasisRTwo}。
% To verify linearly independence we looked at 
% linear combinations of the set's members that total to
% the zero vector 
% $c_1\vec{\beta}_1+c_2\vec{\beta}_2=\binom{0}{0}$.
% The resulting 
% calculation shows that such a combination is unique,
% that $c_1$ must be $0$ and $c_2$ must be $0$.
为了验证该集合张成空间，我们考察其和为空间中
一个成员的线性组合
$c_1\vec{\beta}_1+c_2\vec{\beta}_2=\binom{x}{y}$。
在那个例子中我们只指出这样的组合存在，即对每个 $x,y$
都存在 $c_1,c_2$，但事实上那个计算还表明这个组合是
唯一的：~$c_1$ 必定是 $(y-x)/2$ 而 $c_2$ 必定是 $2x-y$。
 
\begin{theorem} \label{th:BasisIffUniqueRepWRT}
%<*th:BasisIffUniqueRepWRT>
在任何向量空间中，一个子集是基\index{basis}
当且仅当空间中的每个向量都可以按一种且仅一种方式
表示为该子集元素的线性组合。
%</th:BasisIffUniqueRepWRT>
\end{theorem}

\noindent 如果两个线性组合具有相同的加项而只是次序不同，
或者仅在添加或删去形如
“\( 0\cdot\vec{\beta} \)”的项上有差别，
我们就认为它们是相同的。


\begin{proof}
%<*pf:BasisIffUniqueRepWRT0>
一个序列是基，当且仅当其中的向量构成一个张成的且线性无关的集合。
一个子集是张成集，当且仅当空间中的每个向量都可以按至少一种方式
表示为该子集元素的线性组合。
所以我们只需证明：
一个张成的子集是线性无关的，当且仅当空间中的每个向量
都可以按至多一种方式表示为该子集元素的线性组合。
%</pf:BasisIffUniqueRepWRT0>

%<*pf:BasisIffUniqueRepWRT1>
考虑把一个向量表示为子集成员的线性组合的两种方式。
将两个和式重排，并在必要时添加一些 \( 0\cdot\vec{\beta}_i \) 项，
使得两个和式按相同的次序组合相同的 \( \vec{\beta} \)：
\( \vec{v}=\lincombo{c}{\vec{\beta}} \) 与
\( \vec{v}=\lincombo{d}{\vec{\beta}} \)。
于是
\begin{equation*}
   \lincombo{c}{\vec{\beta}}=\lincombo{d}{\vec{\beta}}
\end{equation*}
成立当且仅当
\begin{equation*}
   (c_1-d_1)\vec{\beta}_1+\dots+(c_n-d_n)\vec{\beta}_n=\zero
\end{equation*}
成立。
所以，断言下面方程中每个系数都为零，
就等价于断言对每个 \( i \) 都有 \( c_i=d_i \)，
也就是说，每个向量都可以按唯一的方式表示为
这些 \( \vec{\beta} \) 的线性组合。
%</pf:BasisIffUniqueRepWRT1>
\end{proof}

\begin{definition} \label{def:RepresentingVectors}
%<*df:RepresentingVectors>
在以 $B$ 为基的向量空间中，
\definend{\( \vec{v} \) 关于 \( B \) 的表示}%
\index{representation!of a vector}\index{vector!representation of} 是
用来把 $\vec{v}$ 表示为基向量的线性组合的系数
所构成的列向量：
\begin{equation*}
  \rep{\vec{v}}{B}
  =
  \colvec{c_1 \\ c_2 \\ \vdots \\ c_n}_{B}
\end{equation*}
其中
\( B=\sequence{\vec{\beta}_1,\dots,\vec{\beta}_n} \) 且
\( \vec{v}=\lincombo{c}{\vec{\beta}} \)。
这些 \( c \) 称为
\definend{\( \vec{v} \) 关于 \( B \) 的坐标}%
\index{coordinates!with respect to a basis}。
%</df:RepresentingVectors>
\end{definition}

\begin{example}  \label{ex:RepPolyWRTA Basis}
在 \( \polyspace_3 \) 中，关于基
\( B=\sequence{1,2x,2x^2,2x^3} \)，
\( x+x^2 \) 的表示是
\begin{equation*}
  \rep{x+x^2}{B}=\colvec[r]{0 \\ 1/2 \\ 1/2 \\ 0}_B
\end{equation*}
这是因为 $x+x^2=0\cdot 1+(1/2)\cdot 2x+(1/2)\cdot 2x^2+0\cdot 2x^3$。
而关于另一个基 \( D=\sequence{1+x,1-x,x+x^2,x+x^3} \)，
表示就不同了。
\begin{equation*}
  \rep{x+x^2}{D}=\colvec[r]{0 \\ 0 \\ 1 \\ 0}_D
\end{equation*}
\end{example}

\begin{remark}
\nearbydefinition{def:basis} 要求基是一个序列，
这样我们才能按某个次序写出这些坐标。
\end{remark}

当只涉及一个基时，我们常常省略标明该基的下标。

\begin{example}
在 \( \Re^2 \) 中，为求向量 $\vec{v}=\binom{3}{2}$ 关于基
\begin{equation*}
  B=\sequence{
              \colvec[r]{1 \\ 1},
              \colvec[r]{0 \\ 2} }
\end{equation*}
的坐标，求解
\begin{equation*}
  c_1\colvec[r]{1 \\ 1}
  +c_2\colvec[r]{0 \\ 2}
  =
  \colvec[r]{3 \\ 2}
\end{equation*}
得到 $c_1=3$ 与 $c_2=-1/2$。
\begin{equation*}
  \rep{\vec{v}}{B}=\colvec[r]{3 \\ -1/2}
\end{equation*}
\end{example}

把表示写成一列，推广了熟悉的情况：~在 \( \Re^n \) 中，
关于标准基 \( \stdbasis_n \)，
起点在原点、终点在
\( (v_1,\dots,v_n) \) 的向量就有这样的表示。
\begin{equation*}
  \rep{\colvec{v_1 \\ \vdots \\ v_n}}{\stdbasis_n}
    =
  \colvec{v_1 \\ \vdots \\ v_n}_{\stdbasis_n}
\end{equation*}
下面是一个例子。
\begin{equation*}
  \rep{\colvec{-1 \\ 1}}{\stdbasis_{n}}
  =\colvec{-1 \\ 1}
\end{equation*}

\begin{remark}  \label{rem:RepNotation}
$\rep{\vec{v}}{B}$ 这个记号不是标准的。
最常见的记号是 $[\vec{v}]_B$，但 $\rep{\vec{v}}{B}$
有一个优点：不容易被误解或被忽视。
\end{remark}

说这个列表示该向量，其含义是：
一组向量之间成立某个线性关系，
当且仅当这组表示之间成立该线性关系。

\begin{lemma}  \label{lem:LinearRelRepresented}
%<*lm:LinearRelRepresented>
设 $B$ 是一个有 $n$~个元素的基，则对任意一组向量，
$a_1\vec{v}_1+\cdots+a_k\vec{v}_k=\vec{0}_{V}$ 当且仅当
$a_1\rep{\vec{v}_1}{B}+\cdots+a_k\rep{\vec{v}_k}{B}=\vec{0}_{\Re^n}$。
%</lm:LinearRelRepresented>
\end{lemma}

\begin{proof}
取定一个基 \( B=\sequence{\vec{\beta}_1,\dots,\vec{\beta}_n} \)，
并设
\begin{equation*}
  \rep{\vec{v}_1}{B}=\colvec{c_{1,1} \\ \vdots \\ c_{n,1}}
  \quad\ldots\quad
  \rep{\vec{v}_k}{B}=\colvec{c_{1,k} \\ \vdots \\ c_{n,k}}
\end{equation*}
于是 $\vec{v}_1=c_{1,1}\vec{\beta}_1+\dots+c_{n,1}\vec{\beta}_n$，等等。
那么 $a_1\vec{v}_1+\dots+a_k\vec{v}_k=\vec{0}$ 就等价于下面这些式子。
\begin{align*}
  \vec{0}
  &=a_1\cdot(c_{1,1}\vec{\beta}_1+\dots+c_{n,1}\vec{\beta}_n)
    +\dots+
    a_k\cdot(c_{1,k}\vec{\beta}_1+\dots+c_{n,k}\vec{\beta}_n)  \\
  &=(a_1c_{1,1}+\dots+a_kc_{1,k})\cdot\vec{\beta}_1
    +\dots+
    (a_1c_{n,1}+\dots+a_kc_{n,k})\cdot\vec{\beta}_n 
\end{align*}
显然，若这些系数为零，则下面的方程成立。  
但由于 \( B \) 是一个基，\nearbytheorem{th:BasisIffUniqueRepWRT}
表明：下面的方程成立当且仅当这些系数为零。
所以这个关系等价于
\begin{align*}
    a_1c_{1,1}+\dots+a_kc_{1,k} &=0    \\
                               &\alignedvdots \\ 
    a_1c_{n,1}+\dots+a_kc_{n,k} &=0
\end{align*}
这就是等价地改写成的列向量形式。
\begin{equation*}
  a_1\colvec{c_{1,1} \\ \vdots \\ c_{n,1}}
  +\dots+
  a_k\colvec{c_{1,k} \\ \vdots \\ c_{n,k}}
  =\colvec{0 \\ \vdots \\ 0}
\end{equation*}
注意，不仅一个关系对一组向量成立当且仅当它对另一组成立，
而且这是同一个关系\Dash 这些 \( a_i \) 是相同的。
\end{proof}

\begin{example}
\nearbyexample{ex:RepPolyWRTA Basis} 求出了
\( x+x^2\in\polyspace_3 \) 关于
\( B=\sequence{1,2x,2x^2,2x^3} \) 的表示。
\begin{equation*}
  \rep{x+x^2}{B}=\colvec[r]{0 \\ 1/2 \\ 1/2 \\ 0}_B
\end{equation*}
关系式
\begin{equation*}
  2\cdot(x+x^2)-1\cdot(2x)-2\cdot(x^2) = 0+0x+0x^2+0x^3
\end{equation*}
由下面这个式子表示。
\begin{equation*}
  2\cdot\rep{x+x^2}{B}-\rep{2x}{B}-2\cdot\rep{x^2}{B}
  =2\cdot\colvec{0 \\ 1/2 \\ 1/2 \\ 0}
   -\colvec{0 \\ 1 \\ 0 \\ 0}
   -2\cdot\colvec{0 \\ 0 \\ 1/2 \\ 0}
  =\colvec{0 \\ 0 \\ 0 \\ 0}
\end{equation*}

\end{example}

表示的主要用途要到后面才出现，但这个定义放在这里，
是因为每个向量都可以按唯一的方式表示为基向量的线性组合，
这一事实是基的一个关键性质，
同时也是为了帮助阐明一个观点。
除其他用途外，为了计算坐标，我们将把注意力
限制在基只含有限多个元素的空间上。
这将从下一小节开始。

